\def\pgfsysdriver{pgfsys-dvipdfm.def} % required for forests inside
\documentclass[aspectratio=169
  , xcolor={svgnames}
  , hyperref={ colorlinks,citecolor=Blue
             , linkcolor=DarkRed,urlcolor=DarkBlue}
  , russian
  ]{beamer}
\usetheme{CambridgeUS}
\usefonttheme{professionalfonts}

\usepackage{pgfpages}
\usepackage{bibentry}
\usepackage{cite}
\def\newblock{\hskip .11em plus .33em minus .07em}


\makeatletter
\@ifclassloaded{beamer}{
  % get rid of header navigation bar
  \setbeamertemplate{headline}{}
  % get rid of bottom navigation symbols
  \setbeamertemplate{navigation symbols}{}
  % get rid of footer
  %\setbeamertemplate{footline}{}
}
{}
\makeatother
%%%%%%%%%%%%%%%%%%%%%%%%%%%%%%%%%%%%%%%%%%%%%
\usepackage{fontawesome}
% \newfontfamily{\FA}{Font Awesome 5 Free} % some glyphs missing
\expandafter\def\csname faicon@facebook\endcsname{{\FA\symbol{"F09A}}}
\def\faQuestionSign{{\FA\symbol{"F059}}}
\def\faQuestion{{\FA\symbol{"F128}}}
\def\faExclamation{{\FA\symbol{"F12A}}}
\def\faUploadAlt{{\FA\symbol{"F093}}}
\def\faLemon{{\FA\symbol{"F094}}}
\def\faPhone{{\FA\symbol{"F095}}}
\def\faCheckEmpty{{\FA\symbol{"F096}}}
\def\faBookmarkEmpty{{\FA\symbol{"F097}}}

\newcommand{\faGood}{\textcolor{ForestGreen}{\faThumbsUp}}
\newcommand{\faBad}{\textcolor{red}{\faThumbsODown}}
\newcommand{\faWrong}{\textcolor{red}{\faTimes}}
\newcommand{\faMaybe}{\textcolor{blue}{\faQuestion}}
\newcommand{\faCheckGreen}{\textcolor{ForestGreen}{\faCheck}}
%%%%%%%%%%%%%%%%%%%%%%%%%%%%%%%%%%%%%%%%%%%%%

\usepackage{fontspec}
\usepackage{xunicode}
\usepackage{xltxtra}
\usepackage{xecyr}
\usepackage{hyperref}

\setmainfont[
 Ligatures=TeX,
 Extension=.otf,
 BoldFont=cmunbx,
 ItalicFont=cmunti,
 BoldItalicFont=cmunbi,
% Scale = 1.1
]{cmunrm}
\setsansfont[
 Ligatures=TeX,
 Extension=.otf,
 BoldFont=cmunsx,
 ItalicFont=cmunsi,
%  Scale = 1.2
]{cmunss}
%\setmainfont[Mapping=tex-text]{DejaVu Serif}
%\setsansfont[Mapping=tex-text]{DejaVu Sans}
%\setmonofont{Fira Code}[Contextuals=AlternateOff]
\setmonofont{Fira Code}[Contextuals=Alternate,Scale=0.9]
\newfontfamily{\myfiracode}[Scale=1.5,Contextuals=Alternate]{Fira Code}
%\setmonofont[Scale=0.9,BoldFont={Inconsolata Bold}]{Inconsolata}

\usepackage{polyglossia}
\setmainlanguage{russian}
\setotherlanguage{english}


%\newfontfamily\dejaVuSansMono{DejaVu Sans Mono}
% https://github.com/vjpr/monaco-bold/raw/master/MonacoB/MonacoB.otf
%\newfontfamily\monacoB{MonacoB}
%%%%%%%%%%%%%%%%%%%%%%%%%%%%%%%%%%%%%%%%%%%%%%%5
\usepackage{soul} % for \st that strikes through
\usepackage[normalem]{ulem} % \sout

\usepackage{stmaryrd}
\newcommand{\sem}[1]{\ensuremath{\llbracket #1\rrbracket}}


\usepackage{listings}
%\lstdefinestyle{style1}{
%  language=haskell,
%  numbers=left,
%  stepnumber=1,
%  numbersep=10pt,
%  tabsize=4,
%  showspaces=false,
%  showstringspaces=false
%}
%\lstdefinestyle{hsstyle1}
%{ language=haskell
%%          , basicstyle=\monacoB
%         , deletekeywords={Int,Float,String,List,Void}
%         , breaklines=true
%         , columns=fullflexible
%         , commentstyle=\color{ForestGreen}
%         , escapeinside=§§
%         , escapebegin=\begin{russian}\commentfont
%         , escapeend=\end{russian}
%         , commentstyle=\color{ForestGreen}
%         , escapeinside=§§
%         , escapebegin=\begin{russian}\color{ForestGreen}
%         , escapeend=\end{russian}
%         , mathescape=true
%%          , backgroundcolor = \color{MyBackground}
%}
%
%\newcommand{\inline}[1]{\lstinline{haskell}{#1}}
%\def\hsinline{\mintinline{haskell}}
%\def\inline{\hsinline}
%
%\lstnewenvironment{hslisting} {
%    \lstset { style={hsstyle1} }
%  }
%  {}
%  
%%%%%%%%%%%%%%%%%%%%%%%%%%%%%%%%%%%%%%%%%%%%%%%%%%%%%%%%%%%  
%%\setmainfont[
%% Ligatures=TeX,
%% Extension=.otf,
%% BoldFont=cmunbx,
%% ItalicFont=cmunti,
%% BoldItalicFont=cmunbi,
%%]{cmunrm}
%%% С засечками (для заголовков)
%%\setsansfont[
%% Ligatures=TeX,
%% Extension=.otf,
%% BoldFont=cmunsx,
%% ItalicFont=cmunsi,
%%]{cmunss}
%% \setmonofont[Scale=0.6]{Monaco}
%
%\usefonttheme{professionalfonts}
%\usepackage{times}
\usepackage{tikz}
\usetikzlibrary{cd}
\usepackage{tikz-cd}
\usepackage{caption}
\usepackage{subcaption}

%\renewtheorem{definition}{برهان}[chapter]
%%\DeclareMathOperator{->}{\rightarrow}
%\newcommand\iso{\ensuremath{\cong}}
%\usepackage{verbatim}
%\usepackage{graphicx}
%\usetikzlibrary{arrows,shapes}

%\usepackage{amsmath}
%\usepackage{amsfonts}
\usepackage{scalerel}
\DeclareMathOperator*{\myvee}{\scalerel*{\vee}{\sum}}
\DeclareMathOperator*{\mywedge}{\scalerel*{\wedge}{\sum}}

%
%\usepackage{tabulary}
%
%% sudo aptget install ttf-mscorefonts-installer
%%\setmainfont{Times New Roman}
%%\setsansfont[Mapping=tex-text]{DejaVu Sans}
%
%%\setmonofont[Scale=1.0,
%%    BoldFont=lmmonolt10-bold.otf,
%%    ItalicFont=lmmono10-italic.otf,
%%    BoldItalicFont=lmmonoproplt10-boldoblique.otf
%%]{lmmono9-regular.otf}
%
\usepackage[cache=true]{minted}
\usemintedstyle{perldoc}

\def\hsinline{\mintinline{haskell}}
\def\mlinline{\mintinline{ocaml}}
% color options
\definecolor{YellowGreen} {HTML}{B5C28C}
\definecolor{ForestGreen} {HTML}{009B55}
\definecolor{MyBackground}{HTML}{F0EDAA}



\institute{матмех СПбГУ}

\addtobeamertemplate{title page}{}{
  \begin{center}{\tiny Дата сборки: \today}\end{center}
}


\lstdefinelanguage{ocamllambda}{
keywords={catch, switch, default, case, with, failwith, exit, true, false, ::},
sensitive=true,
commentstyle=\small\itshape\ttfamily,
keywordstyle=\ttfamily\textbf,
identifierstyle=\ttfamily,
basewidth={0.5em,0.5em},
columns=fixed,
mathescape=true,
fontadjust=true,
literate={->}{{$\to$}}3 {===}{{$\equiv$}}1  {Scru}{{$\bullet$}}2,
morecomment=[s]{(*}{*)}
}

\usepackage{subcaption}
\usepackage{etoolbox}


\usepackage{exercise}
\usepackage{tikz}
\usepackage{tikz-qtree}
\usetikzlibrary{trees}
\usepackage[edges]{forest}
\forestset{.style={
%  for tree={l=1em, l sep=1em, s sep=1em}
  forked edges,
    for tree={    grow'=0,    draw,    align=c,    font=\sffamily,
        rounded corners  }
  }}

\newcommand{\lstquot}[1]{``\lstinline{#1}''}
\newcommand{\sembr}[1]{\llbracket{#1}\rrbracket}
%\newcommand\false{$f\!alse$}
%\newcommand\myif{i\!f}


\def\transarrow{\xrightarrow}
\newcommand{\setarrow}[1]{\def\transarrow{#1}}

\def\padding{\phantom{X}}
\newcommand{\setpadding}[1]{\def\padding{#1}}

\def\subarrow{}
\newcommand{\setsubarrow}[1]{\def\subarrow{#1}}

\newcommand{\trule}[2]{\dfrac{#1}{#2}}
\newcommand{\crule}[3]{\dfrac{#1}{#2},\;{#3}}
\newcommand{\withenv}[2]{{#1}\vdash{#2}}
\newcommand{\trans}[3]{{#1}\transarrow{\padding{\textstyle #2}\padding}\subarrow{#3}}
\newcommand{\ctrans}[4]{{#1}\transarrow{\padding#2\padding}\subarrow{#3},\;{#4}}
\newcommand{\llang}[1]{\mbox{\lstinline[mathescape]|#1|}}
\newcommand{\pair}[2]{\inbr{{#1}\mid{#2}}}
\newcommand{\inbr}[1]{\left<{#1}\right>}
\newcommand{\highlight}[1]{\color{red}{#1}}
%\newcommand{\ruleno}[1]{\eqno[\scriptsize\textsc{#1}]}
\newcommand{\ruleno}[1]{\mbox{[\textsc{#1}]}}
\newcommand{\rulename}[1]{\textsc{#1}}
\newcommand{\inmath}[1]{\mbox{$#1$}}
%\newcommand{\lfp}[1]{fix_{#1}}
%\newcommand{\gfp}[1]{Fix_{#1}}
\newcommand{\vsep}{\vspace{-2mm}}
\newcommand{\supp}[1]{\scriptsize{#1}}
\renewcommand{\sembr}[1]{\llbracket{#1}\rrbracket}
\newcommand{\cd}[1]{\texttt{#1}}
%\newcommand{\free}[1]{\boxed{#1}}
%\newcommand{\binds}{\;\mapsto\;}
%\newcommand{\dbi}[1]{\mbox{\bf{#1}}}
%\newcommand{\sv}[1]{\mbox{\textbf{#1}}}
%\newcommand{\bnd}[2]{{#1}\mkern-9mu\binds\mkern-9mu{#2}}
%\newcommand{\meta}[1]{{\mathcal{#1}}}
%\newcommand{\dom}[1]{\mathtt{dom}\;{#1}}
%\newcommand{\primi}[2]{\mathbf{#1}\;{#2}}
%\renewcommand{\dom}[1]{\mathcal{D}om\,({#1})}
%\newcommand{\ran}[1]{\mathcal{VR}an\,({#1})}
%\newcommand{\fv}[1]{\mathcal{FV}\,({#1})}
%\newcommand{\tr}[1]{\mathcal{T}r_{#1}}
\newcommand{\diseq}{\not\equiv}
%\newcommand{\reprfunset}{\mathcal{R}}
%\newcommand{\reprfun}{\mathfrak{f}}
%\newcommand{\cstore}{\Omega}
%\newcommand{\cstoreinit}{\cstore_\epsilon^{init}}
%\newcommand{\csadd}[3]{add(#1, #2 \diseq #3)}  %{#1 + [#2 \diseq #3]}
%\newcommand{\csupdate}[2]{update(#1, #2)}  %{#1 \cdot #2}
\newcommand{\primi}[1]{\ensuremath{\mathbf{#1}}}
\newcommand{\sem}[1]{\llbracket #1 \rrbracket}
\newcommand{\ir}{\ensuremath{\mathcal{S}}}
\usepackage{tikz}
\newcommand*\circled[1]{\tikz[baseline=(char.base)]{
    \node[shape=circle,draw,inner sep=1pt] (char) {#1};}}

\let\emptyset\varnothing
\let\eps\varepsilon

% for fancy table
%\newcommand{\lheadl}[2]{\multicolumn{1}{|>{\centering\arraybackslash}m{#1}|}{{#2}}}
%\newcommand{\head}[2]{\multicolumn{1}{>{\centering\arraybackslash}m{#1}}{\textbf{\small #2}}}
%\newcommand{\headll}[2]{\multicolumn{1}{>{\centering\arraybackslash}m{#1}||}{\textbf{\small #2}}}
%\newcommand{\lheadll}[2]{\multicolumn{1}{|>{\centering\arraybackslash}m{#1}||}{\textbf{\small #2}}}
%\newcommand{\headl}[2]{\multicolumn{1}{>{\centering\arraybackslash}m{#1}|}{\textbf{\small #2}}}
%\usepackage{longtable}
%\newcommand{\nodata}{}
%\newcommand{\tablenotemark}[1]{#1}

\newcommand{\contributions}{
\begin{enumerate}
\item[I] Спроектировали синтез с помощью комбинации \emph{реляционных интерпретаторов} на \miniKanren{}
\item[II] Заменили $\forall$ на \emph{конечный} набор примеров
\item[III] Сделали оптимизацию методом ветвей и границ с помощью нового примитива \miniKanren{}: \emph{ограничение на структуру (structural constraint)}
\end{enumerate}
}

\title[Компиляция сопоставления с образцом]{Компиляция сопоставления с образцом с помощью синтеза программ}
%\subtitle{Relational Synthesis for Pattern Matching}

\date{19 декабря 2020}
\author{Косарев Дмитрий} 
\institute[]{%\normalfont
Лаборатория языковых инстументов \\
Миниконференция 2020}


\newcommand{\verbatimfont}[1]{\def\verbatim@font{#1}}
\usepackage{verbatimbox}

\forestset{
    .style={
        for tree={
            base=bottom,
            child anchor=north,
            align=center,
            s sep+=1cm,
    straight edge/.style={
        edge path={\noexpand\path[\forestoption{edge},thick,-{Latex}] 
        (!u.parent anchor) -- (.child anchor);}
    },
    if n children={0}
        {tier=word, draw, thick, rectangle}
        {draw, diamond, thick, aspect=2},
    if n=1{%
        edge path={\noexpand\path[\forestoption{edge},thick,-{Latex}] 
        (!u.parent anchor) -| (.child anchor) node[pos=.2, above] {Y};}
        }{
        edge path={\noexpand\path[\forestoption{edge},thick,-{Latex}] 
        (!u.parent anchor) -| (.child anchor) node[pos=.2, above] {N};}
        }
        }
    }
}
  
%\AtBeginSection[]
%{
%  \begin{frame}<beamer>
%    \frametitle{Оглавление}
%    \tableofcontents[currentsection,currentsubsection]
%  \end{frame}
%} 
%\AtBeginSubsection[]
%{
%  \begin{frame}<beamer>
%    \frametitle{Оглавление}
%    \tableofcontents[currentsection,currentsubsection]
%  \end{frame}
%}

\begin{document}
{
\begin{frame}[fragile]
  \begin{tabular}{p{5.5cm} p{5.5cm}}
   \begin{center}
%      \includegraphics[height=1.5cm]{pictures/jetbrainsResearch.pdf}
    \end{center}
    &
    \begin{center}
   %   \includegraphics[height=1.5cm]{pictures/SPbGU_Logo.png}
    \end{center}
  \end{tabular}
  \titlepage
\end{frame}
}

%\maketitle

% For every picture that defines or uses external nodes, you'll have to
% apply the 'remember picture' style. To avoid some typing, we'll apply
% the style to all pictures.
\tikzstyle{every picture}+=[remember picture] 

% By default all math in TikZ nodes are set in inline mode. Change this to
% displaystyle so that we don't get small fractions.
\everymath{\displaystyle}
%\begin{comment}
% Uncomment these lines for an automatically generated outline.
\begin{frame}{Сопоставление с образцом}
Существенная часть функционального программирования
\vspace{1cm}

%\includegraphics{pictures/SPbGU_Logo.png}

Два основных подхода:
\begin{itemize}
\item диаграммы решений
\begin{itemize}
\item минимизируется количество проверок
\end{itemize}
\item автомат с возвратами
\begin{itemize}
\item минимизируется размер кода
\end{itemize}%\pause
\item \textbf{синтез программ} (наш подход)
\end{itemize}
\end{frame}

\section{Обзор}

\defverbatim[colored]{\matchA}{
\begin{lstlisting}
match x,y,z with
| _,F,T -> 1
| F,T,_ -> 2
| _,_,F -> 3
| _,_,T -> 4
\end{lstlisting}
}

\defverbatim[colored]{\matchB}{
\begin{lstlisting}
if x then
  if y then
    if z then 4 else 3
  else
    if z then 1 else 3
else
  if y then 2
  else
    if z then 1 else 3
\end{lstlisting}
}

\defverbatim[colored]{\matchC}{
\begin{lstlisting}
if y then
  if x then
    if z then 4 else 3
  else 2
else
  if z then 1 else 3
\end{lstlisting}
}


\defverbatim[colored]{\forestB}{
\begin{forest} 
[\texttt{x}%, tikz={\draw[{Latex}-, thick] (.north) --++ (0,1);}
    [\texttt{y}
          [\texttt{z}
                [4] 
                [3] 
            ]   
         [\texttt{z}
               [1] 
               [3] 
           ]    
    ]   
    [\texttt{y}
        [2] 
        [\texttt{z}
            [1] 
            [3] 
        ]   
    ]   
] 
\end{forest}
}

\defverbatim[colored]{\forestC}{
\begin{forest} 
[\texttt{y}%, tikz={\draw[{Latex}-, thick] (.north) --++ (0,1);}
    [\texttt{x}
          [\texttt{z}
                [4] 
                [3] 
            ]   
         [2]    
    ]
    [\texttt{z}
               [1] 
               [3] 
           ]
]
\end{forest}
}


\begin{frame}[fragile]{Пример: использование диаграмм решений \only<1>{(1/2)}\only<2>{(2/2)}}
%\begin{figure}[ht]
%\begin{subfigure}[t]{.49\textwidth}
%
%\end{subfigure}
%\begin{subfigure}[t]{.49\textwidth}
%
%\end{subfigure}
%\end{figure}

\begin{figure}[ht]
\begin{subfigure}[t]{0.25\linewidth}
\begin{minipage}{7cm}
\matchA
\end{minipage}
\caption{Пример}
\end{subfigure}
\hspace{0.5cm}
\begin{subfigure}[t]{0.32\linewidth}
\only<1>{\begin{minipage}{5cm}
\matchB
\caption{Полуоптимальная схема компиляции}
\end{minipage}}
\only<2>{\begin{minipage}{5cm}
\matchC
\caption{Оптимальная схема компиляции}
\end{minipage}}
\end{subfigure}
\hspace{0.5cm}
\begin{subfigure}[t]{0.3\linewidth}
\only<1>{\begin{minipage}{7cm}
\forestB
\end{minipage}}
\only<2>{\begin{minipage}{7cm}
\forestC
\end{minipage}}
\caption{Диаграмма решений}
\end{subfigure}
\caption{Компиляция сопоставления с образцом бывает нетривиальной (пример из~\cite{maranget2008}).}
%\label{fig:match-example}
\end{figure}
\end{frame}

\begin{comment}


\defverbatim[colored]{\btCompAone}{
\begin{lstlisting}[language=ocamllambda]


    (switch lx with case []: 1
      default:     )
               
      (switch ly with case []: 2
        default:     )
                  
      (switch lx with
        case (::):
          (switch ly with
            case (::) : 3
            default:     )
        default:     )   
                                
\end{lstlisting}
}

\defverbatim[colored]{\btCompA}{
\begin{lstlisting}[language=ocamllambda]
catch
  (catch
    (switch lx with case []: 1
      default: exit)
    with (catch
      (switch ly with case []: 2
        default: exit)
    with (catch
      (switch lx with
        case (::):
          (switch ly with
            case (::) : 3
            default: exit)
        default: exit))))
with (failwith "Partial match")
\end{lstlisting}
}

\defverbatim[colored]{\btCompBone}{
\begin{lstlisting}[language=ocamllambda]
      
         
    (switch lx with
      case []: 1
      case (::) :
        (switch ly with
          case (::): 3
          default:     ))
      
    (switch ly with
       case []: 2
       default:     )
                               
\end{lstlisting}
}

\defverbatim[colored]{\btCompB}{
\begin{lstlisting}[language=ocamllambda]
catch
  (catch
    (switch lx with
      case []: 1
      case (::) :
        (switch ly with
          case (::): 3
          default: exit))
  with
    (switch ly with
       case []: 2
       default: exit)
with (failwith "Partial match")
\end{lstlisting}
}

\begin{frame}[fragile]{Пример: Компиляция в автомат с возвратами}
\begin{minipage}{0.45\linewidth}\vspace{0pt}
\[
(P\rightarrow L)=    \begin{pmatrix}
      \texttt{[]} & \texttt{\_} & \rightarrow & 1 \\
      \texttt{\_} & \texttt{[]} & \rightarrow & 2  \\
      \texttt{x::xs} & \texttt{y::ys} & \rightarrow & 3
    \end{pmatrix}
\]
\vspace{1em}
\uncover<2->{
\[
(P_1\rightarrow L_1)=    \begin{pmatrix}
      \texttt{[]} & \texttt{\_} & \rightarrow & 1 \\
    \end{pmatrix}
\]
\[
(P_2\rightarrow L_2)=    \begin{pmatrix}
       \texttt{\_} & \texttt{[]} & \rightarrow & 2  \\
    \end{pmatrix}
\]
\[
(P_3\rightarrow L_3)=    \begin{pmatrix}
      \texttt{x::xs} & \texttt{y::ys} & \rightarrow & 3
    \end{pmatrix}
\]}
\end{minipage}\hspace{.3cm}
\begin{minipage}{0.45\linewidth}%\vspace{-1em}
\begin{minipage}{0.45\linewidth}
\only<3>{\btCompAone}
\only<4>{\btCompA}
%\vspace{1em}
\end{minipage}
\end{minipage}
\end{frame}

\begin{frame}[fragile]{Пример: Компиляция в автомат с возвратами + оптимизация}
\begin{minipage}[t]{0.45\linewidth}\vspace{0pt}
\uncover<1>{
\[
(P\rightarrow L)=    \begin{pmatrix}
      \texttt{[]} & \texttt{\_} & \rightarrow & 1 \\
      \texttt{\_} & \texttt{[]} & \rightarrow & 2  \\
      \texttt{x::xs} & \texttt{y::ys} & \rightarrow & 3
    \end{pmatrix}
\]}
%\vspace{1em}
\uncover<2->{
\[
(P\rightarrow L)=    \begin{pmatrix}
      \texttt{[]} & \texttt{\_} & \rightarrow & 1 \\
      \texttt{x::xs} & \texttt{y::ys} & \rightarrow & 3\\
      \texttt{\_} & \texttt{[]} & \rightarrow & 2  
    \end{pmatrix}
\]
}\vspace{1cm}
\uncover<3->{
\[
(P_1\rightarrow L_1)=    \begin{pmatrix}
      \texttt{[]} & \texttt{\_} & \rightarrow & 1 \\
      \texttt{x::xs} & \texttt{y::ys} & \rightarrow & 3
    \end{pmatrix}
\]
\[
(P_2\rightarrow L_2)=    \begin{pmatrix}
       \texttt{\_} & \texttt{[]} & \rightarrow & 2  \\
    \end{pmatrix}
\]
}
\end{minipage}\hspace{.5cm}
\begin{minipage}[t]{0.45\linewidth}\vspace{0pt}
\begin{minipage}{0.44\linewidth}
\only<1-2>{\vspace{18em}}
\only<4>{\btCompBone\vspace{1em}}
\only<5>{\btCompB}
\end{minipage}
\end{minipage}
\end{frame}


\begin{frame}{Ещё оптимизации для автоматов с возвратами}
\begin{itemize}
\item Использование информации о полноте
\item Оптимизация конструкций \lstinline[language=ocamllambda]{exit}, новый синтаксис \lstinline[language=ocamllambda]{exit N}
\begin{itemize}
\item Дешёвая поддержка охранных выражений (pattern guards)
\end{itemize}
\item Ещё кое-что, для избегания повторных проверок
\end{itemize}
\vspace{2em}
В компиляторе \OCaml{}  используется  реализация из~\cite{maranget2001}.
\end{frame}

%%%%%%%%%%%%%%%%%%%%%%%%%%%%%%%%%%%%%%%%%%%%%%%%%%%%%%%%%%%%%%%%%%%%%
\defverbatim[colored]{\haskellDemoHead}{
\begin{lstlisting}[language=haskell]
f :: Bool -> Bool
f True = True 
f x    = f x 

y :: Bool
y = f x 
\end{lstlisting}
}

\defverbatim[colored]{\haskellDemo}{
\begin{lstlisting}[language=haskell]
case x,y of 
  True, False -> 1
  True, True  -> 2
  False,False -> 3
  False,True  -> 4
\end{lstlisting}
}

\begin{frame}[fragile]{Языки с ленивой семантикой}
\begin{minipage}[t]{0.45\linewidth}
\begin{minipage}{0.5\linewidth}
\uncover<2->{\haskellDemoHead}

\haskellDemo
\end{minipage}
\end{minipage}\hspace{0.1cm}
\begin{minipage}[t]{0.45\linewidth}
\uncover<3>{
\begin{minipage}{1\linewidth}
В \textit{ленивых} языках у нас нет свободы в выборе первого столбца $\Rightarrow$ в \textit{строгих} больше простора для оптимизаций
\end{minipage}}
\end{minipage}
\end{frame}

\begin{frame}{Понятие "хорошей" скомпилированной программы}
\begin{minipage}[t]{0.45\linewidth}
В литературе упоминаются~\cite{Scott2000WhenDM} следующие эвристики
\begin{itemize}
\item Слева-направо
\item Справа-налево
\item Small branching factor (малый коэффициент ветвления)
\item Large branching factor 
\item Leaf edges 
\item Arity factor
\item Artificial rule
\item и другие
\end{itemize}
\end{minipage}\hspace{1cm}
\begin{minipage}[t]{0.45\linewidth}
В подавляющем большинстве случаев формально более оптимальные программы показывают сходную производительность
\end{minipage}
\end{frame}
\end{comment}


%%%%%%%%%%%%%%%%%%%%%%%%%%%%%%%%%%%%%%%%%%%%%%%%%%%%%%%%%%%%%%%
\section{Заслуги работы}

\begin{frame}{Заслуги работы <<Реляционный синтез сопоставления с образцом>>}
\Large
%"Реляционный синтез сопоставления с образцом"
\contributions
\end{frame}

\subsection{Синтез с помощью реляционных интерпретаторов}

\begin{frame}{Реляционные интерпретаторы}
Семейство языков \miniKanren{} -- реинкарнация логического программирования\\

Программы представляются как вычислимые отношения ("реляции", англ. relation)

\[
  \texttt{Interpret}: \texttt{Program} \times \texttt{Input} \times \texttt{Result}
\]

Их можно запускать\footnote{Если интерпретатор написан достаточно аккуратно} в разные стороны:
\[
  \texttt{Interpret}_{\rightarrow}: \texttt{Program} \times \texttt{Input} \rightarrow \texttt{Result}
\]
\[
  \texttt{Synthesize}_{\leftarrow}:   \texttt{Input} \times \texttt{Result} \rightarrow \texttt{Program}
\]
\end{frame}

\begin{frame}[fragile]{Синтаксис двух языков сопоставления с образцом}
\begin{figure}[ht]
\begin{subfigure}[t]{0.4\linewidth}
\uncover<1->{
\[
match^o: \mathcal{V}\times \mathcal{P}^* \times  \mathbb{N}
\]}
Язык сопоставления с образцом
\[
 \begin{array}{rcll}
    \mathcal{C} & = & \{ C_1^{k_1}, \dots, C_n^{k_n} \}\\
    \mathcal{V} & = & \mathcal{C}\,\mathcal{V}^*\\  
    \mathcal{P} & = & \_ \mid \mathcal{C}\,\mathcal{P}^*
 \end{array}
\]

\[
 \begin{array}{c}
\trans{\inbr{v;\,p_1,\dots,p_k}}{}{i}\\
1\leqslant i\leqslant k+1
 \end{array}
\]
\end{subfigure}
\hspace{0.5cm}
\begin{subfigure}[t]{0.5\linewidth}
\uncover<1->{
\[
eval^o_{\mathcal S}: \mathcal{V}\times \mathcal{S} \times  \mathbb{N}
\]}
Язык скомпилированного представления
\[
\begin{array}{rcl}
  \mathcal M & = & \bullet \\
  &   & \mathcal M\,[\mathbb{N}] \\
  \ir & = & \primi{return}\,\mathbb{N} \\
  &   & \primi{switch}\;\mathcal{M}\;\primi{with}\; [\mathcal{C}\; \primi{\rightarrow}\; \ir]^*\;\primi{otherwise}\;\ir
\end{array}
\]
\end{subfigure}
\end{figure}
\vspace{2em}
Опущены для простоты: типы, охранные выражения, переменные в паттернах и т.д.

\end{frame}

\begin{frame}{Алгоритм синтеза}
%High-level description:
\[
%depth^o\,v\,n \;\; \wedge \;\;
\forall v \quad \forall  (1\leqslant\!i\leqslant \!k\!+\!1) \quad
%\mbox{\lstinline|fresh ($i$)|}\; \{
(match^o\,v\,\,p_1,\dots,p_k\,\,i) \;\; \Leftrightarrow \;
eval^o_{\mathcal S}\,v\,\circled{?}\,i
%\}
\]

\begin{itemize}
%\item $depth^o\,v\,n$ generates examples bound by depth $n$
%\item $\forall v$ -- all possible scrutinees
%\item $\forall i$ -- all possible branches
\item $match^o\,v\,\,p_1,\dots,p_k\,\,i$ -- интерпретатор языка сопоставления с образцом, для каждого сопоставляемого выражения (scrutinee) $v$ выдает номер ветви $i$ 
\item $\circled{?}$ -- программа, которую надо синтезировать
\item $eval^o_{\mathcal S}\,v\,\circled{?}\,i$ -- a интерпретатор скомпилированного представления $\mathcal S$, которая проверяет, что синтезированная программа $\circled{?}$ на примерах $v$ выдает правильные номера ветвей  $i$
\end{itemize}
\vspace{1cm}
\begin{itemize}
\item[\faGood] Интерпретаторы  $match^o$ и $eval^o$ легко реализовать (и для расширений сопоставления с образцом тоже должно быть легко)
\item[\faBad] \miniKanren{} с disequality constraints не умеет работать с кванторами $\forall$
\end{itemize}
\end{frame}

\subsection{Создание конечного набора примеров}

\begin{frame}{Избавление от $\forall$. Создание конечного набора примеров}
Для каждого сопоставления с образцом  мы знаем:
\begin{itemize}
\item тип сопоставляемого выражения
\item все образцы, которые используются
\begin{itemize}
\item можем посчитать максимальную глубину образцов
\end{itemize}
\end{itemize}

\begin{alertblock}{\textbf{Идея}}
Переберём всех жителей типа сопоставляемого выражения до некоторой глубины, и будем использовать этих жителей как примеры
\end{alertblock}
\vspace{1cm}

\begin{itemize}
%\item[\faGood] Is correct (proved in Coq)
\item[\faBad] В худшем случае ---~экспоненциальное количество примеров
\end{itemize}
\end{frame}

\begin{comment}
\begin{frame}{Our Method: Synthesis + Relational Interpreter}
Do \emph{not compile} with specific algorithm but \emph{synthesize} compiled representation on large enough but finite set of examples
\vspace{1cm}

Fast $\Leftrightarrow$ Small: synthesized program should contain less checks.
\vspace{1cm}

We are using relational programming, more precisely OCanren~\cite{OCanrenWeb} from miniKanren~\cite{MiniKanrenWeb} family.
\vspace{1cm}


Our repo on Github: ~\cite{Repo}.
\end{frame}
\end{comment}


\defverbatim[colored]{\exampleNotAFullSet}{
\begin{lstlisting}{ocaml}
match (s : unit list) with 
| _ :: _ :: _ -> 1
| _           -> 2
\end{lstlisting}
}

\begin{comment}
\begin{frame}{Пример 1: Полный набор примеров из трёх штук}
\centering
\begin{minipage}{.5\textwidth}
\begin{minipage}{.5\textwidth}
\exampleNotAFullSet
\end{minipage}\end{minipage}

\begin{figure}
\begin{subfigure}[t]{0.3\linewidth}
\begin{subfigure}[b][4cm][t]{0.2\linewidth}
\begin{tikzpicture}[sibling distance=3em,
  every node/.style = {shape=rectangle, rounded corners, minimum height=.5cm,
    draw, align=center, font=\ttfamily
    ,top color=white, bottom color=blue!20}]]
  \node {Cons}
    child { node {\phantom{1}\_\phantom{1}} }
    child { node {Cons}
      child { node  {\phantom{1}\_\phantom{1}} }
      child { node {\phantom{1}\_\phantom{1}} } 
    };
\end{tikzpicture}
\end{subfigure}
\hspace{2cm}
\begin{subfigure}[b][4cm][t]{0.1\linewidth}
\begin{tikzpicture}[sibling distance=3em,
  every node/.style = {shape=rectangle, rounded corners, minimum height=.5cm,
    draw, align=center, font=\ttfamily
    ,top color=white, bottom color=blue!20}]]
  \node {\phantom{1}\_\phantom{1}};
\end{tikzpicture}
\end{subfigure}
\caption{Два образца}
\end{subfigure}
\hspace{1.5cm}
\begin{subfigure}[t]{0.5\linewidth}
\begin{subfigure}[b][4cm][t]{0.2\linewidth}
\begin{tikzpicture}[sibling distance=3em,
  every node/.style = {shape=rectangle, rounded corners, minimum height=.5cm,
    draw, align=center, font=\ttfamily
    ,top color=white, bottom color=green!20}]]
  \node {Cons}
    child { node {()} }
    child { node {Cons}
      child { node {()} }
      child { node { Nil } } 
    };
\end{tikzpicture}
\end{subfigure}
\hspace{1cm}
\begin{subfigure}[b][4cm][t]{0.15\linewidth}
\begin{tikzpicture}[sibling distance=3em,
  every node/.style = {shape=rectangle, rounded corners, minimum height=.5cm,
    draw, align=center, font=\ttfamily
    ,top color=white, bottom color=green!20}]]
  \node {Cons}
    child { node {()} }
    child { node {Nil} } ;
\end{tikzpicture}
\end{subfigure}
\hspace{1cm}
\begin{subfigure}[b][4cm][t]{0.15\linewidth}
\begin{tikzpicture}[sibling distance=3em,
  every node/.style = {shape=rectangle, rounded corners, minimum height=.5cm,
    draw, align=center, font=\ttfamily
    , top color=white, bottom color=green!20}]]
  \node {Nil};
\end{tikzpicture}
\end{subfigure}
\caption{Три примера}
\end{subfigure}
\end{figure}
\end{frame}


\begin{frame}[fragile]{Пример 2: Недостаточно полный набор примеров}
\begin{minipage}{0.4\linewidth}
\begin{lstlisting}{ocaml}
match (s : unit list) with 
| []  -> 1
| _   -> 2
\end{lstlisting}
\end{minipage}
\begin{minipage}{0.55\linewidth}
Два образца:\\
\begin{tikzpicture}[sibling distance=3em,
  every node/.style = {shape=rectangle, rounded corners, minimum height=.5cm,
    draw, align=center, font=\ttfamily
    ,top color=white, bottom color=blue!20}]]
  \node { Nil };
\end{tikzpicture}
\begin{tikzpicture}[sibling distance=3em,
  every node/.style = {shape=rectangle, rounded corners, minimum height=.5cm,
    draw, align=center, font=\ttfamily
    ,top color=white, bottom color=blue!20}]]
  \node {\phantom{1}\_\phantom{1}};
\end{tikzpicture}
\vspace{1em}

Набор примеров, ограниченных глубиной 1, состоит только из одного примера\\
\begin{tikzpicture}[sibling distance=2em,
  every node/.style = {shape=rectangle, rounded corners, minimum height=.5cm,
    draw, align=center, font=\ttfamily
    , top color=white, bottom color=green!20}]]
  \node {Nil};
\end{tikzpicture}
\vspace{1em}

Данная программа ведет себя согласовано на наборе примеров \\
\[
 \begin{array}{ll}
    \primi{switch} \dots \primi{with}   & \\
    |\ \text{Nil} \rightarrow 1  & \\  
    |\ \primi{otherwise} \rightarrow  1 & 
 \end{array}
\]

но, очевидно, неправильная
\end{minipage}
\end{frame}


\begin{frame}{Текущий алгоритм для получения примеров}
\begin{figure}
\begin{subfigure}[b]{0.75\linewidth}
Если кратко:
\begin{itemize}
\item Вычислить глубину образцов $h$
\item Синтезируем всех жителей, но
\item на глубине $h+1$ используем заранее подготовленного жителя соответствующего типа
\end{itemize}
\vspace{1cm}

\begin{itemize}
\item[\faBad] В худшем случае ---~экспоненциальное количество примеров
\end{itemize}
\end{subfigure}
\end{figure}
\end{frame}
\end{comment}


\subsection{Оптимизации}

\begin{comment}



\begin{frame}[fragile]{Оптимизация: откидывание эквивалентных программ}
Очевидно, что реляционный интерпретатор языка \ir{} может перебирать различные эквивалентные программы
\begin{center}
  \begin{minipage}[t]{0.2\linewidth}
  \begin{center}
  \begin{lstlisting}[language=ocamllambda,gobble=2]
  switch Scru with 
  true -> ...
  false -> ...
  \end{lstlisting}
  \end{center}
  \end{minipage}\hspace{.5cm}
  \begin{minipage}[t]{0.2\linewidth}
  \huge
  \begin{center}
  \[
  \Leftrightarrow
  \]
  \end{center}
  \end{minipage}\hspace{.5cm}
  \begin{minipage}[t]{0.2\linewidth}
  \begin{lstlisting}[language=ocamllambda,gobble=2]
  switch Scru with 
  false -> ...
  true -> ...
  \end{lstlisting}
  \end{minipage}
\end{center}
С этой проблемой удалось побороться
\begin{itemize}
\item Задав порядок на ветвях \primi{switch}ей, используя информацию о типе
\item Это несколько <<сломало>> интерпретатор, но для синтеза это не существенно
\end{itemize}

\end{frame}

\defverbatim[colored]{\improvementA}{
\begin{lstlisting}{ocaml}
match (s : bool * bool * bool) with 
| (_,_,F) -> 1
| (_,_,T) -> 2
\end{lstlisting}
}

\begin{frame}{Оптимизация: сокращение необходимого набора примеров}
Пример:
\improvementA
\vspace{1em}

Мы можем во время компиляции обнаружить, что
\begin{itemize}
\item Всего $2^3$ жителей у типа \lstinline=bool * bool * bool=
\item Не надо проверять, что сопоставляемое выражение --- это тройка
\item Не надо заглядывать в 1ю и 2ю компоненты, так как важна только третья 
\end{itemize}
\vspace{1em}


Итого, может запускать синтез только на двух примерах $\{(\mathcal{B},\mathcal{B},\text{T}),(\mathcal{B},\mathcal{B},\text{F})\}$
(где $\mathcal{B}$ ---~это любое значение типа \lstinline=bool=)
если мы \textbf{запретим "заглядывание" в ненужные поддеревья} примеров
\end{frame}
\end{comment}

\begin{frame}{Оптимизация методом ветвей и границ}
Текущий результат синтеза хранится в  $\circled{?}$ и во время поиска 
мы \emph{только уточняем} результат, добавляя новые конструкции $\primi{switch}$ \\

\begin{alertblock}{\textbf{Идея}}
Если текущее приближение ответа длиннее, чем уже найденный ответ --- прерырываем поиск в этой ветке
\end{alertblock}
\vspace{2em}

%Требует модификации примитива \miniKanren{} \lstinline=run=: для каждого найденного ответа
%\begin{itemize}
%\item считаем размер
%\item обновляем минимальный найденный ответ
%\end{itemize}

Отсечение ветвей поиска делается помощью нового примитива --- \emph{ограничения на структуру (structural constraint)}

\end{frame}









%%%%%%%%%%%%%%%%%%%%%%%%%%%%%%%%%%%%%%%%%%%%%%%%
%\section{Реализация}
%
%\begin{frame}{Реализация}
%Мы пользуемся \OCanren{} --- типобезопасным встраиванием \miniKanren{} в \OCaml{}.
%\vspace{2em}
%
%В процессе используется \noCanren{}~\cite{RelConversion} для порождения кода на \OCanren{}~\cite{OCanren,OCanrenWeb}
%\vspace{1em}
%
%Основная часть (два реляционных интерпретатора + порождение примеров) сделаны с помощью \noCanren{}.
%\vspace{1em}
%
%Репозиторий проекта~\cite{Repo}
%\end{frame}

%%%%%%%%%%%%%%%%%%%%%%%%%%%%%%%%%%%%%%%%%%%%%%%%
\section{Проблемы}

\defverbatim[colored]{\synthInputA}{
\begin{lstlisting}
match $\bullet$ with
(_, false, true) -> 1
(false, true, _) -> 2
(_, _, false) -> 3
(_, _, true) -> 4
\end{lstlisting}
}
\defverbatim[colored]{\synthResultA}{
\begin{lstlisting}[language=ocamllambda]
switch $\bullet$[0] with  
| true -> 
    (switch $\bullet$[1] with  
    | true -> 
        (switch $\bullet$[2] with true -> 4 | _ -> 3)
    | _ -> 
        (switch $\bullet$[2] with true -> 1 | _ -> 3))  
| _ -> 
    (switch $\bullet$[1] with  
    | true -> 2   
    | _ -> 
         (switch $\bullet$[2] with true -> 1 | _ -> 3))
\end{lstlisting}
}

\defverbatim[colored]{\synthResultB}{
\begin{lstlisting}[language=ocamllambda]
switch Scru[0] with  
| true -> 
    (switch Scru[2] with  
    | true -> 
       (switch Scru[1] with true -> 4 | _ -> 1)
    | _ -> 3 )  
| _ -> 
    (switch Scru[1] with  
    | true -> 2   
    | _ -> (switch Scru[2] with  true -> 1 | _ -> 3))
\end{lstlisting}
}

\defverbatim[colored]{\synthResultC}{
\begin{lstlisting}[language=ocamllambda]
switch $\bullet$[1] with  
| true -> 
    (switch $\bullet$[0] with  
    | true -> 
        (switch $\bullet$[2] with true -> 4 | _ -> 3)
    | _ -> 2) 
| _ -> 
    (switch $\bullet$[2] with true -> 1 | _ -> 3)
\end{lstlisting}
}

\begin{frame}[fragile]{Пример синтезированной программы: сопоставление выражения типа \lstinline[language=ocaml]{bool*bool*bool}}
\begin{minipage}[c][7cm][t]{0.35\linewidth}\vspace{0em}
\begin{minipage}{0.35\linewidth}
\synthInputA
\end{minipage}
\end{minipage}
\begin{minipage}{0.63\linewidth}\vspace{0em}
\begin{onlyenv}<1>\begin{minipage}{0.63\linewidth}
Ответ размера  6 (за 1.6 секунд)
\synthResultA
\end{minipage}\end{onlyenv}
%
\begin{onlyenv}<2>\begin{minipage}{0.63\linewidth}
Ответ размера 5 (за +0.4 секунд)
\synthResultB
\end{minipage}\end{onlyenv}
\begin{onlyenv}<3>\begin{minipage}{0.63\linewidth}
Ответ размера 4 (за +0.7 секунд)
\synthResultC
\end{minipage}\end{onlyenv}
\end{minipage}
\end{frame}

%%%%%%%%%%%%%%%%%%%%%%%%%%%%%%%%%%%%%%%%%%%%%%%%

\begin{frame}[fragile]{Проблемы с производительностью}
\begin{figure}
\begin{subfigure}[t]{0.5\linewidth}
\begin{lstlisting}[basicstyle=\small,language=ocaml]
match a,s,c with
| (_,_,Ldi i::_) -> 1
| (_,_,Push::_)  -> 2
| (Int _,Val (Int _)::_,IOp _::_) -> 3
| (Int _,_,Test (_,_)::c) -> 4
| (Int _,_,Test (_,_)::c) -> 5
| (_,_,Extend::_) -> 6
| (_,_,Search _::_) -> 7
| (_,_,Pushenv::_) -> 8
| (_,Env e::s,Popenv::_) -> 9
| (_,_,Mkclos cc::_) -> 10
| (_,_,Mkclosrec _::_) -> 11
| (Clo (_,_), Val _::_, Apply::_) -> 12
| (_,(Code _::Env _::_),[]) -> 13
| (_,[],[]) -> 14
\end{lstlisting}
\end{subfigure}
\hspace{1cm}
\begin{subfigure}[t]{0.35\linewidth}\vspace{0em}
Интерпретатор PCF (mini-ML) из статьи Г.Плоткина, 1977
\vspace{4em}

Сейчас не работает, потому что слишком много (11102) примеров
\begin{itemize}
\item большая глубина образцов
\item много конструкторов в типах
\end{itemize}
%Doesn't currently work because the types are too large (11102 examples  generated).\\

%But for a reduced example we can synthesize the answer
%\vspace{1cm}
\end{subfigure}
\end{figure}
\end{frame}

\begin{comment}

\begin{frame}[fragile]{Проблемы с производительностью (2/2)}
\begin{figure}
\begin{subfigure}[b][7cm][b]{0.35\linewidth}\vspace{0em}
\begin{lstlisting}[basicstyle=\small]
type code = 
| Push 
| Ldi of int 
| IOp of int 
| Int of int 
type prog = code list 
type item = 
| Val of code 
| Env of int 
| Code of int
type stack = item list 

match (code,stack,prog) with
| (_, _, (Ldi _)::_)  -> 1
| (_, _, (Push _)::_) -> 2
\end{lstlisting}
\begin{onlyenv}<1>
\vspace{4mm}
\end{onlyenv}
\begin{onlyenv}<2>
\begin{lstlisting}[basicstyle=\small,aboveskip=-0.5em]
| (Int _, _, (IOp _)::_) -> 3
\end{lstlisting}
\end{onlyenv}
\end{subfigure}
\hspace{.5cm}
\begin{subfigure}[b]{0.65\linewidth}
\begin{overlayarea}{8cm}{6cm}
Сокращённый пример 
\begin{itemize}
\item по типам 
\item по веткам
\end{itemize}
Для двух веток  надо 5 примеров\\\vspace{1em}

\begin{onlyenv}<2->
Для трёх веток и тех же типов уже необходимо 20 примеров
\begin{itemize}
\item за 1,5с получим 1й ответ размера 7
\item ещё через полсекунды --- 2й и 3й (последний) ответы размера 6 и 5, соответственно
\item в конце оно тратит 10с, чтобы доказать, что более коротких ответов не существует
\end{itemize}
%\begin{lstlisting}[basicstyle=\small,aboveskip=-0.5em]
%| (Int _, _, (IOp _)::_) -> 3
%\end{lstlisting}
\end{onlyenv}
\end{overlayarea}
\vspace{1cm}
\end{subfigure}
\end{figure}

\end{frame}
\end{comment}


\section{Заключение}

\begin{frame}{Результаты}% и задачи на будущее}
Достижения: \contributions
\vspace{1em}

На маленьких примерах подход работает корректно, ... но на больших есть проблемы с производительностью 
\vspace{1em}
\end{frame}

\begin{frame}{Доклады и публикации}
\begin{enumerate}
\item TEASE-LP 2020 (доклад)
\item MiniKanren Workshop при ICFP 2020 (доклад)
\item APLAS-2020 (доклад и публикация~\cite{matchingPublication})
\end{enumerate}
\vspace{4em}
\begin{center}
{\Huge Конец}
\end{center}
\end{frame}

%\begin{frame}{Пути дальнейшего улучшения}
%\begin{itemize}
%\item Построение входного логического значения для  ограничений на структуру можно делать эффективнее (ленивые вычисления)
%\item Использование конечнодоменных ограничений вместо disequality constraints 
%\item Мемоизация сейчас никак не используется, т.к.  disequality constraints
%\item Создание меньшего числа примеров
%\item Другое представление языка \ir{} с использованием конструкций \primi{exit}
%\begin{itemize}
%\item тут может потребоваться номинальная унификация
%\end{itemize}
%\end{itemize}
%\end{frame}
%
%\begin{frame}
%\begin{center}
%{\Huge Спасибо!}
%\end{center}
%\end{frame}


\begin{frame}%[t, allowframebreaks]
\frametitle{Литература}
\bibliographystyle{amsalpha}
\bibliography{references}
\vspace{1cm}
\end{frame}


\begin{frame}[fragile,noframenumbering]{Критерий минимизации для синтезированных программ}
\begin{figure}
\begin{subfigure}[b]{0.3\linewidth}
\[
  \begin{array}{l}
  \primi{switch}\;\mathcal{M}\;\primi{with}\; \\
  \mathcal{C}_1\; \primi{\rightarrow}\; \ir_1\\
  \dots \\
  \mathcal{C}_n\; \primi{\rightarrow}\; \ir_n\\
  \primi{otherwise}\;\ir\\
  \end{array}\\
\]
\[
  \Updownarrow
\]
\[
  \begin{array}{l}
  \primi{if}\; \mathcal{M} = \mathcal{C}_1\;\primi{then}\; \ir_1\\
  \dots \\
  \primi{else\ if}\; \mathcal{M} = \mathcal{C}_n\; \primi{then}\;  \ir_n\\
  \primi{else}\;\ir\\
  \end{array}\\
\]
\vspace{1cm}

\end{subfigure}
\hspace{1cm}
\begin{subfigure}[b]{0.6\linewidth}
Интуиция: один \primi{switch} с $n$ случаев можно примерно закодировать в  $n$ \primi{if}ов\\

Будем считать, что размер
\begin{itemize}
\item \primi{switch} --- это число веток
\item \primi{return} равен 0
\item программы целиком --- сумма размеров всех входящих в неё \primi{switch}
\end{itemize}
Наш критерий минимизации: уменьшение размера синтезированной программы \\

Но могут быть другие: глубина, коэффициент ветвления, и т.д.
\end{subfigure}
\end{figure}

\end{frame}


\begin{frame}[noframenumbering]{Ограничение на структуру (structural constraint)}
Новый примитив:
\begin{itemize}
\item Принимает промежуточное представление значений и конвертирует их (в текущем состоянии) до логических %a logic value and performs reification in a current state
\item Принимает предикат для логических значений %Takes a predicate which inspects the reified value 
\item Если результат слишком большой ---~прекратить поиск% If there are too many branches, stops the search 
($failure^o$)
\item Иначе продолжать поиск, не меняя состояние 
%Otherwise continues search without changing the state 
($success^o$)
\end{itemize}
\vspace{1em}
%Peculiarities:
Особенности:
\begin{itemize}
\item Используется для вычисления размера текущего решения
\item Может учитывать или не учитывать disequality constraints
\item Можно использовать, чтобы реализовать $absent^o$ и подобные играничения
\item Работает с логическим представлением (reified) значений $\Rightarrow$ медленно
\end{itemize}
\end{frame}

\end{document}
